\section{The homological carrier}
\label{sec:carrier}

\subsection{Where the flat band comes from}

Project the Kogut--Susskind Hamiltonian \eqref{eq:KS} onto the
fundamental, charge-odd, one-plaquette manifold. In three spatial
dimensions a fundamental plaquette excitation carries one of three face
orientations, and at second order neighbouring faces communicate
through a shared link. The projected problem is therefore a
three-orbital Bloch Hamiltonian, and it looks at first like a
tight-binding model.

It is not an engineered one. The signs of the matrix elements are fixed
by the orientation of the plaquette boundary and by charge conjugation,
and non-abelian group integration generates, exactly, the oriented
edge--face boundary Gram operator
\begin{equation}
  B(k)B(k)^{\dagger},
  \qquad B = \partial_2^{\dagger},
\end{equation}
with the coefficient of \S\ref{sec:second} in front of it. The
structure of the result is a clean separation:
\begin{equation}
  \boxed{\ \text{colour dynamics} \longrightarrow t_N,
  \qquad
  \text{cellular geometry} \longrightarrow BB^{\dagger}. \ }
  \label{eq:separation}
\end{equation}
Representation theory fixes the amplitude; the cubical chain complex
fixes the kernel and the spectrum. Neither borrows from the other.

\subsection{The projected spectrum}

\begin{theorem}[Second-order projected spectrum]
\label{thm:projspec}
The projected second-order spectrum consists of one
momentum-independent level together with a doubly degenerate branch
shifted by $t_N u^2 q(k)$, where
\begin{equation}
  q(k) \;=\; 4\sum_{j=1}^{3} \sin^2\!\big(k_j/2\big).
\end{equation}
\end{theorem}

The flat level is exact and it is exact for a reason: the flat carrier
is $\ker \partial_2$.

\begin{theorem}[The carrier and its multiplicity]
\label{thm:multiplicity}
Elementary cube boundaries are compact localized states of the carrier;
three wrapping sheets complete it on the torus. For $L \ge 3$ the exact
multiplicity is
\begin{equation}
  (L^3 - 1) + 3 \;=\; L^3 + 2 .
\end{equation}
\end{theorem}
\gid{RESULT:UNIVERSAL\_CELLULAR\_HODGE\_SPECTRUM}

The two contributions are structurally different and the distinction
matters later. The $L^3 - 1$ cube boundaries are local and finitely
supported. The three sheets are not: they wrap the torus, and they are
the entire zero-momentum content of the carrier. \S\ref{sec:wilson}
returns to the consequence --- that no finitely supported exactly-flat
local operator can see the $\Gamma$ point at all.

\subsection{Third order, and the coefficients that survive}

For $\SU(3)$, exact third-order perturbation theory stays inside the
same boundary-factorized algebra. It does not generate a new operator
structure; it renormalizes the coefficient in front of the same
$BB^{\dagger}$:
\begin{equation}
  H_{\mathrm{eff},-}
  \;=\; E_{\mathrm{flat}}(u)\, I
  \;+\; \left(\frac{5}{612}u^2 + \frac{1975}{124848}u^3\right) BB^{\dagger}
  \;+\; O(u^4),
  \label{eq:thirdorder}
\end{equation}
with
\begin{equation}
  E_{\mathrm{flat}}(u)
  \;=\; \frac{8}{3} + u + \frac{11}{306}u^2
        - \frac{109151}{249696}u^3 .
\end{equation}

That \eqref{eq:thirdorder} remains boundary-factorized through third
order is the non-trivial part. It is what makes the fourth order the
first place where the carrier can couple to its complement, and
\S\ref{sec:fourth} is about exactly that coupling.

\subsection{What is and is not novel here}

This program does not claim that incidence flatness, cellular compact
localized states, or homological counting of zero modes are new. All
three have a substantial prior literature --- line-graph constructions,
compact localized states, the necessity of non-contractible states on
periodic systems, rectangular factorizations and index arguments, and
closely related closed-surface structures in two-form spin liquids ---
and a recent independent construction treats face-graph and higher-cell
flat bands directly.

What is claimed is that the flatness here is \emph{dynamical}: it is
produced by exact strong-coupling matrix elements of non-abelian
$\SU(N)$ Hamiltonian lattice gauge theory, with the amplitude fixed by
group representation theory and the kernel fixed by the cubical chain
complex, as in \eqref{eq:separation}. The flat band is not put in; it
comes out.

\begin{scopenote}
Theorem~\ref{thm:projspec} and \eqref{eq:thirdorder} are finite-volume,
finite-order statements about a projected flux sector. They are not a
continuum glueball theorem and not a mass-gap theorem, and the source
manuscript says so in its own abstract.
\end{scopenote}
